% ============================================================================
% Appendix: the sense-making question sets and the corpus source.
% Two sets of 60 questions each, both produced by GraphRAG's own question-generation protocol (5 potential users x 3
% tasks x 4 questions); the generator's only input is a textual description of the corpus -- it sees no index,
% community partition, or any method's answers, so it is structurally unbiased toward the evaluated methods. The difference between the two sets comes
% from two independent samples at temperature=0.7 happening to fall into different question styles: the first set starts with how/what
% (asking for cross-community synthesis), the second with identify/determine (asking for locating and enumerating).
% This difference systematically changes the middle of the ranking (see the discussion of Table~\ref{tab:arena-appendix}).
% ============================================================================
\section{Question sets and corpus provenance}
\label{app:questions}

\stitle{Corpus.} All end-to-end experiments use MultiHop-RAG~\cite{tang2024multihoprag},
the news collection GraphRAG itself evaluates on. It contains $609$ English news
articles published between 2023-09-26 and 2023-12-25 and drawn from $49$ outlets,
the largest being Sporting News ($101$ articles), TechCrunch ($97$), The Verge
($45$), Polygon ($44$) and The Guardian ($29$). By category the corpus is
sports ($211$), technology ($172$), entertainment ($114$), business ($81$), science ($21$) and health ($10$).
We index the first $304$ articles in publication order and ingest the next $23$
as one update batch, giving the $327$-document index of Section~\ref{sec:exp:e2e}.

\stitle{Question generation.} GraphRAG ships no question set with reference
answers; its own evaluation protocol instead asks a language model to imagine
potential users of a corpus, tasks each would perform, and questions each task
raises, given only a short prose description of the corpus. We follow that
protocol verbatim with $5$ users $\times$ $3$ tasks $\times$ $4$ questions, run
twice at temperature $0.7$ to obtain two independent sets of $60$. The generator
sees only the corpus description---never an index, a community partition, or any
system's answer---so it cannot favour a method by construction. What the two runs
did not hold fixed is interrogative style, and that turned out to matter: the
first set opens with \emph{how} ($26$), \emph{what} ($16$) and \emph{can} ($15$),
the second with \emph{identify} ($12$), \emph{determine} ($8$), \emph{analyze}
($6$) and \emph{evaluate} ($5$). Median question length is $96$ and $88$
characters respectively.

\subsection{Set A: synthesis questions (\texorpdfstring{$q_0$--$q_{59}$}{q0-q59})}
\begin{footnotesize}
\begin{enumerate}\setlength{\itemsep}{1pt}
\item What are the most frequently mentioned entities across all categories from September to December 2023?
\item How has the frequency of articles on specific topics changed from September to December 2023?
\item Can you detect any patterns in the type of articles (e.g., opinion pieces vs. news reports) over the four months?
\item Which categories saw the highest increase in article count from September to December 2023?
\item What is the overall sentiment towards major corporations mentioned in the business category?
\item How does public sentiment towards a particular celebrity change from September to December 2023?
\item Can you determine if there is a correlation between positive mentions of a technology company and their stock performance?
\item What is the sentiment towards environmental issues as reported in the health and science category?
\item How do the most common themes in the entertainment category compare with those in the technology category?
\item What are the key differences in the types of entities discussed in sports versus business articles?
\item Can you identify any significant overlaps in the entities discussed in health and science articles compared to technology articles?
\item How does the frequency of articles about international events differ between the sports and business categories?
\item How consistent is the use of sources across different categories?
\item What is the distribution of original reporting versus rehashing of press releases or other media?
\item Can you assess the diversity of viewpoints presented in the articles?
\item How often do articles reference the same entities, and what does this imply about the scope of coverage?
\item How did major global events (e.g., political elections, natural disasters) influence the content of the articles across categories?
\item What are the top three most covered events or issues in each category, and how did these topics evolve over the four months?
\item Can you identify any long-term trends in the coverage of certain events that may indicate a shift in editorial focus?
\item How does the coverage of local events compare with that of global events in the sports and technology categories?
\item What themes have emerged consistently across different categories throughout the period?
\item How do the themes in the entertainment category compare with those in the health and science category?
\item Can you identify any thematic shifts in the coverage of business news from September to December 2023?
\item How do the thematic focuses of the sports category differ from those in the technology category?
\item What are the most popular topics among readers based on the frequency of mentions?
\item How do consumer interests vary between different categories?
\item What emerging trends can be identified from the data that could inform marketing strategies?
\item Can you determine which entities are most likely to drive consumer interest in the business category?
\item Which brands have seen the highest increase in visibility across the four months?
\item How does brand visibility differ across the various categories?
\item Can you identify any brands that have maintained consistent visibility throughout the period?
\item Which brands are most associated with positive sentiment in the technology category?
\item How does the coverage of a specific industry (e.g., healthcare, tech) affect its visibility and perception?
\item What are the most influential companies in terms of driving industry-wide discussions?
\item Can you determine if there is a correlation between increased coverage and changes in stock prices for listed companies?
\item How does the coverage of a particular industry compare with its actual market performance?
\item What are the most frequently discussed health-related issues during the period?
\item How do public health concerns in the news compare with those in medical journals?
\item Can you identify any trends in the coverage of new health technologies or treatments?
\item How does the coverage of mental health issues compare with that of physical health issues?
\item What is the public's sentiment towards specific health issues (e.g., obesity, diabetes)?
\item How does the coverage of a particular health issue affect public awareness and behavior?
\item Can you determine if there is a correlation between the volume of articles and public health outcomes?
\item How does the coverage of a specific health issue (e.g., vaccination campaigns) vary across different regions?
\item What are the most discussed technological advancements in healthcare during the period?
\item How do these technological advancements compare with their actual adoption rates in clinical practice?
\item Can you identify any trends in the integration of technology into healthcare practices?
\item How does the coverage of telemedicine compare with traditional healthcare methods?
\item What are the most frequently discussed technological advancements during the period?
\item How do these technological advancements compare with the latest innovations in the tech industry?
\item Can you identify any emerging trends in the tech sector that could present new opportunities?
\item How does the coverage of a specific technology (e.g., AI, blockchain) compare with its actual adoption rate?
\item What are the most discussed tech companies during the period, and how do they compare with each other?
\item How does the coverage of a specific tech company (e.g., Apple, Google) affect its market position?
\item Can you determine if there is a correlation between increased coverage and market share growth for specific tech companies?
\item How does the coverage of a particular tech trend (e.g., cybersecurity) vary across different industries?
\item What is the public's sentiment towards specific technologies (e.g., smartphones, smart homes)?
\item How does the coverage of a particular technology (e.g., electric vehicles) affect consumer perception and behavior?
\item Can you identify any trends in the adoption of new technologies based on the coverage?
\item How does the coverage of a specific technology (e.g., virtual reality) compare with its actual usage rates?
\end{enumerate}
\end{footnotesize}

\subsection{Set B: enumeration questions (\texorpdfstring{$q_{60}$--$q_{119}$}{q60-q119})}
\begin{footnotesize}
\begin{enumerate}\setlength{\itemsep}{1pt}\setcounter{enumi}{60}
\item How has the frequency of coverage on specific topics changed from September to December?
\item What are the most common themes across different categories of news articles?
\item Which entities have been covered consistently across multiple categories, and why might this be significant?
\item Identify any emerging or declining trends in the types of stories being reported.
\item Can you identify any biases in the reporting based on the tone and language used?
\item How does the coverage of major events differ between categories?
\item Are there any discrepancies in the portrayal of certain entities across different news articles?
\item Examine if there is a correlation between the source of information and the angle of reporting.
\item How have global events influenced the local news coverage during the period?
\item Identify any news articles that were directly affected by international developments.
\item Compare the impact of domestic versus international news events on the local audience.
\item Determine if there are any patterns in how different categories of news articles respond to the same global event.
\item Calculate the overall sentiment score for each category of news articles over the period.
\item Identify any significant shifts in sentiment scores for specific entities over time.
\item Compare the sentiment scores of different entities within the same category to see if there are any notable differences.
\item Analyze the relationship between the volume of news articles and the overall sentiment score.
\item Using clustering algorithms, group similar articles together based on their content and topics.
\item Determine the number of clusters that best represent the diversity of the dataset.
\item Evaluate the coherence of the clusters formed, ensuring they align with expected categories.
\item Identify any outliers or anomalies within the clusters and explain their presence.
\item Based on the historical data, predict which topics are likely to gain more attention in the next quarter.
\item Assess the predictive power of certain features (e.g., entity mentions, article length) in forecasting future trends.
\item Evaluate the accuracy of your predictions using a validation set of recent articles.
\item Discuss any challenges in predicting future trends based on the current dataset.
\item Which demographic segments are most likely to engage with the news articles in each category?
\item Analyze the content to determine which topics are most appealing to different age groups.
\item Identify any gaps in the current coverage that could attract new audiences.
\item Determine the effectiveness of the current targeting strategies based on the engagement levels.
\item Which types of articles tend to perform better on social media platforms?
\item Analyze the engagement metrics to identify the most effective headlines and formats.
\item Determine the optimal posting times for maximizing reach and engagement.
\item Evaluate the impact of multimedia elements (e.g., images, videos) on article performance.
\item Based on the dataset, what content themes should be prioritized in future articles to maintain high engagement?
\item Identify any underrepresented topics that could be included to diversify the content portfolio.
\item Assess the balance between different categories to ensure a well-rounded content strategy.
\item Develop a plan for incorporating trending topics into the content calendar.
\item Track changes in public opinion on specific issues over the period based on the articles.
\item Identify any shifts in public sentiment towards certain entities or topics.
\item Analyze the role of influential figures in shaping public opinion through the articles.
\item Determine the effectiveness of news articles in influencing public opinion on key issues.
\item Examine the extent to which news articles have influenced policy decisions during the period.
\item Identify any instances where news coverage led to policy changes or public debates.
\item Assess the role of media in highlighting or downplaying certain issues in the policy agenda.
\item Discuss the implications of media influence on policy-making processes.
\item Analyze the relationship between the news coverage and societal trends during the period.
\item Identify any correlations between news articles and changes in societal attitudes or behaviors.
\item Explore the impact of news articles on public discourse and cultural norms.
\item Discuss the broader implications of the findings for understanding the role of media in society.
\item Review the consistency and quality of reporting across different categories.
\item Identify any areas where the reporting standards could be improved or standardized.
\item Assess the accuracy and reliability of the information presented in the articles.
\item Determine the effectiveness of fact-checking processes in maintaining editorial integrity.
\item Analyze the factors that contribute to higher reader engagement in different categories.
\item Identify strategies to increase the reach and impact of news articles.
\item Determine the most effective ways to integrate multimedia elements into the articles.
\item Evaluate the impact of interactive features (e.g., polls, quizzes) on reader engagement.
\item Assess the balance between different types of content (e.g., opinion pieces, investigative reports) in the current mix.
\item Identify any gaps in the current content mix that could be filled to better serve the audience.
\item Determine the most effective ways to incorporate user-generated content into the news articles.
\item Evaluate the potential for partnerships with other media outlets to enhance the content mix.
\end{enumerate}
\end{footnotesize}
